\section{The Method}

Every claim in this paper is one experiment. It builds the exact substrate, runs the law, measures one thing, and returns a verdict. The whole simulator is open, and every experiment in this paper can be rebuilt, stepped, and rerun from it \cite{pollard2026vibetest}. The substrate is discrete and deterministic, so there are no real numbers and no randomness inside it. Where a trend is wanted, it is found by growing the run rather than by averaging over random draws, and real numbers appear only as measured outputs, never in the base.

Each verdict carries a depth grade. The grade keeps a passing test from being mistaken for a deep one, and it says how much was put in by hand.

\begin{table}[ht]
\centering
\begin{tabular}{@{}cp{0.66\textwidth}@{}}
\toprule
grade & what it establishes \\
\midrule
L3 & an emergent and novel result, with a control that could have failed \\
L2 & a known physical construction reproduced on the substrate \\
L1 & a known mathematical fact reproduced on the substrate \\
L0 & circular, the answer was put in by hand, kept only as a consistency check \\
\bottomrule
\end{tabular}
\caption{The depth rubric. Only L3 is both new and guarded.}
\label{tab:rubric}
\end{table}

Behind the grade is a single rule. A result counts only when its control could have come out the other way. A finding means something precisely because there was a way for it to fail. One more rule keeps the program from cheating.

\begin{principle}[Report the negative]
If a phenomenon does not emerge from the eight elements, that is recorded as a negative, never patched by adding a ninth thing.
\end{principle}

One more constraint governs the whole construction from below.

\begin{principle}[Discrete at the base]
The foundation is finite. There are no real numbers, no continuum, and no smooth manifolds at the bottom, only integers and trits a machine can hold exactly. Continuous mathematics enters solely as the large-scale limit of the discrete law, with the coarse-graining made explicit, so the smoothness of physics is a thing to be explained rather than assumed.
\end{principle}
